\section{Norm-One Mass and the Geometric Transfer}
\label{geo:section}

Let $K/F$ be quadratic with involution $\iota$, let $K$ be totally
imaginary, and let $F$ have signature $(b,c)$, where $b>0$. We keep these
assumptions in Sections~\ref{fw:section} and~\ref{pf:section}.
Write
\[
 d=[F:\mathbb Q]=b+2c,\qquad \theta=c/d,
 \qquad |\Delta_K|=\lambda^{2d},\quad
 |\Delta_F|=\lambda^d,\quad \ell=\log\lambda.
\]
Assume that $K/F$ is unramified at all finite places. The discriminants
are field discriminants, and archimedean integrals use Lebesgue measure
on unweighted complex coordinates. We count ordered edges and divide
by two to obtain the stated bounds for unordered edges.

We use Dirichlet's unit theorem, Kronecker's theorem, trace duality,
Hilbert's theorem~90, and the analytic class-number formula
\cite{MilneANT,MilneCFT}.

\subsection{The Unit Mass}

Put $E_K=\mathcal O_K^\times$, $E_F=\mathcal O_F^\times$,
$U=\ker(N:E_K\to E_F)$, and $I_U=[E_F:N(E_K)]$.
At each complex place of $F$, choose its two extensions to $K$ and set
\[
 l_j(x)=\frac12\bigl(\log|\sigma_j^+(x)|-
                         \log|\sigma_j^-(x)|\bigr).
\]
On $U$, this is $\log|\sigma_j^+(x)|$. Its image is a full lattice in
$\mathbb R^c$ of covolume $R_U$, where $R_U=1$ if $c=0$.
Let $w_K$ be the number of roots of unity in $K$,
$h_{\rm rel}=h_K/h_F$, $\kappa=|\ker(\mathrm{Cl}(F)\to\mathrm{Cl}(K))|$,
and $L_F(s)=\zeta_K(s)/\zeta_F(s)$. The following identities give the
mass of the norm-one units needed in the lattice construction.

\begin{proposition}\label{geo:mass}
We have that
\[
 \kappa=I_U/2^{b-1},\qquad
 R_K/R_F=2^{c-1}I_U R_U,
\]
and consequently
\begin{equation}\label{geo:mass-density}
 \frac{w_K}{h_{\rm rel}\kappa R_U}
 =\frac{2^{d-1}\pi^{b+c}}{\lambda^{d/2}L_F(1)}.
\end{equation}
Here $R_K,R_F$ are ordinary regulators, using twice the logarithm of
absolute value at a complex place.
\end{proposition}

\begin{proof}
Dirichlet's theorem gives invariant and anti-invariant rational unit
ranks $b+c-1$ and $c$. The Herbrand quotient is multiplicative in exact
sequences, equals one for finite groups, and is $2$ or $1/2$ on a trivial
or sign copy of $\mathbb Z$. Applying this to the finite-index sum of
the two eigenspace lattices gives the quotient $2^{b-1}$, without
requiring the integral unit lattice to split.
Its degree-zero Tate group has size $I_U$, so
$|H^1(\langle\iota\rangle,E_K)|=I_U/2^{b-1}$.

Capitulation injects into this $H^1$, since a relation
$\mathfrak a\mathcal O_K=(x)$ gives the cocycle $x/\iota x$ and a
coboundary gives an invariant generator. Conversely, Hilbert 90
writes every norm-one unit as $x/\iota x$. The fractional ideal $(x)$ is
invariant. It descends from $F$, since exponents agree on split prime
pairs and every finite ramification index is one. This proves
surjectivity and the formula for $\kappa$.

The norm map on rational unit spaces is onto the invariant space, since
the norm of an $F$-unit is its square. Its kernel has rank $c$.
The logarithmic kernel in $U$ consists of roots of unity by Kronecker's
theorem. Every root of unity of $K$ has norm one, since its norm is a
torsion unit of $F$ and is positive at a real place.
Thus the logarithmic kernel has exactly $w_K$ elements.

For the regulator identity, delete a coordinate over a real place of
$F$ in both weighted logarithmic unit spaces. The norm adds the two
coordinates above a complex place of $F$. Replacing a pair $(x,y)$ by
$(x,x+y)$ has determinant one. On the kernel its first coordinate is
$2l_j$, giving covolume $2^cR_U$. The logarithmic norm image has index
$I_U/2$ in the logarithmic $F$-unit lattice. The torsion index is two,
since $-1$ is not a norm at a complexifying real place. Multiplication
of the kernel and image covolumes gives
$R_K/R_F=2^cR_U(I_U/2)$.

The analytic class-number formula
\cite[Chapter~V, Theorem~2.4]{MilneCFT}, with $w_F=2$, gives
\[
 h_{\rm rel}\frac{R_K}{R_F}
 =\frac{w_K\lambda^{d/2}L_F(1)}{2^{c+1}\pi^{b+c}}.
\]
Substitution proves \eqref{geo:mass-density}.
\end{proof}

Choose finitely many primes $\mathfrak p$ of $F$ split as
$\mathfrak P\iota\mathfrak P$ in $K$, and integers $k_{\mathfrak p}\ge0$.
Set $P_0=\prod_{\mathfrak p}(k_{\mathfrak p}+1)$. The one-sided ideals
$B_{\boldsymbol j}=\prod\mathfrak P^{j_{\mathfrak p}}$ have a fiber of
size at least $P_0/(\kappa h_{\rm rel})$ in
$\mathrm{Cl}(K)/\mathrm{im}\,\mathrm{Cl}(F)$. If
$B_{\boldsymbol j}/B_{\boldsymbol j_0}=(x)\mathfrak a\mathcal O_K$, then
\[
 J_{\boldsymbol j}/J_{\boldsymbol j_0}=(x/\iota x),\qquad
 J_{\boldsymbol j}=\prod\mathfrak P^{j_{\mathfrak p}}
                (\iota\mathfrak P)^{k_{\mathfrak p}-j_{\mathfrak p}}.
\]
Thus $I=J_{\boldsymbol j_0}^{-1}$ contains this many disjoint cosets
$\beta_jU$ of norm-one elements. Their principal ideals distinguish the
cosets, and
\[
 N_K(I)=\prod_{\mathfrak p}(N_F\mathfrak p)^{-k_{\mathfrak p}}.
\]
For uniform rational-prime data define
\begin{equation}\label{geo:HJ}
 H=\sum_p\frac{k_p\log p}{e_p},\qquad
 J=\sum_p\frac{\log(k_p+1)}{e_pf_p}.
\end{equation}
Here $e_p,f_p$ are the absolute indices at the selected primes, and all
$F$-primes above them split in $K/F$. Then $N_K(I)=e^{-dH}$ and
$P_0=e^{dJ}$. Prime-ideal mixtures give the corresponding sums with a
bounded rounding error, provided the finite list of exponents is fixed.

\subsection{Unit Averaging and Translation}

For $h\in\mathbb R^c$, rescale each complex pair by $(e^{-h_j},e^{h_j})$,
leaving the $b$ compact coordinates unchanged. Denote the resulting ideal
lattice by $\Lambda_h$. Its determinant, independently of $h$, is
\begin{equation}\label{geo:determinant}
 D=2^{-d}\lambda^d e^{-dH}.
\end{equation}
For any nonnegative integrable $\Phi$ on $\mathbb R^c$, averaging over a
fundamental parallelepiped $P_U$ of $l(U)$ gives exactly
\begin{equation}\label{geo:unfold}
 \frac1{R_U}\int_{P_U}
 \sum_{\beta\in\bigcup_j\beta_jU}\Phi(l(\beta)-h)\,dh
 =\frac{q w_K}{R_U}\int_{\mathbb R^c}\Phi(u)\,du,
 \qquad q\ge\frac{P_0}{\kappa h_{\rm rel}}.
\end{equation}
This follows from tiling and Tonelli's theorem for the full unit lattice.
We use logarithmic Haar measure $du$ and angular measures of total mass
one.

Let $\Omega$ be bounded and invariant under rotations in each complex
coordinate. Let $V$ be its volume and let $\Phi(u)$ be its overlap with
a displacement having compact moduli one and pair moduli $(e^{u_j},e^{-u_j})$.
This function is nonnegative, bounded, and compactly supported in $u$.
Parameterize additive translates by $B_hr$, $r\in[0,1)^{2d}$, where
$B_h$ is a continuous basis matrix of $\Lambda_h$. The point count $N(h,r)$
and inherited ordered-edge count $E(h,r)$ are measurable countable sums
of indicators. Averaging $r$ by lattice tiling gives
\[
 \mathbb E_r N(h,r)=V/D,\qquad
 \mathbb E_r E(h,r)=D^{-1}\sum_\beta\Phi(l(\beta)-h).
\]
If a uniform point bound $N\le N_{\max}$ holds for every $h,r$, all these
quantities are integrable. Averaging also in $h$, and weighting parameter
space by $N$, selects one pair $(h,r)$ with $N>0$ and
\begin{equation}\label{geo:joint-average}
 \frac{E(h,r)}{N(h,r)}\ge
 \frac{q w_K}{R_U}\frac{\int\Phi}{V}.
\end{equation}
The vertex bound and the edge bound hold for the same rescaling and
translate.

\subsection{Uniform Lattice Counting}

The real Euclidean dual of the unweighted complex ideal lattice is the
coordinatewise complex conjugate of the image of
$2\mathfrak D_K^{-1}I^{-1}$. Indeed its real pairing is the trace pairing
divided by two. Reciprocal deformations invert on the dual. Consequently
every nonzero dual vector satisfies
\begin{equation}\label{geo:dual-product}
 \prod_{i=1}^d|\xi_i|\ge\mu^d,
 \qquad \mu=2e^{H/2-\ell},\qquad
 \sum_i|\xi_i|\ge d\mu.
\end{equation}
The product assertion follows by taking the square root of the absolute
norm and using $|N_K(x)|\ge1/(|\Delta_K|N_K(I))$ for
$0\ne x\in\mathfrak D_K^{-1}I^{-1}$.
Here conjugation acts coordinatewise and need not agree with $\iota$.

Suppose a smooth positive weight $F$ on $\mathbb C^d$, formed as a product
over the chosen one- and two-coordinate blocks, has normalized
Fourier bound, for the convention $e^{-2\pi i x\cdot\xi}$,
\[
 |\widehat F(\xi)|/\int F\le M^d
                  \exp\left(-\sigma\sum_i|\xi_i|\right),
 \qquad M\ge1,\quad \sigma>0.
\]
Put $R=d\mu$ and use the norm $\|\xi\|_*=\sum_i|\xi_i|$.
Every two distinct dual points are separated by at least $R$, so open
balls of radius $R/2$ centered at them are disjoint. Comparing volumes
inside the ball of radius $(j+3/2)R$ bounds the number of points in
$jR\le\|\xi\|_*<(j+1)R$ by
$(2j+3)^{2d}\le5^{2dj}$. Thus the sufficient condition
\begin{equation}\label{geo:poisson-condition}
 \sigma\mu\ge\log M+2\log5+2
\end{equation}
makes the nonzero normalized Fourier sum at most
$\sum_{j\ge1}e^{-2dj}<1$.
For Gaussian and polynomial Student profiles, we obtain
\begin{equation}\label{geo:pointwise-poisson}
 \sum_{x\in\Lambda_h+r}F(x)\le2\int F/D,
\end{equation}
uniformly in $h,r$. To justify Poisson summation, multiply $F$ by
$e^{-\varepsilon|x|^2}$. The result is Schwartz. Convolution with the
normalized positive Fourier Gaussian multiplies the Fourier envelope by
its exponential norm moment, which tends to one in each fixed dimension.
The strict tail bound persists for sufficiently small $\varepsilon$.
Schwartz Poisson summation and monotone convergence on the primal side
prove \eqref{geo:pointwise-poisson}. This limit is taken before $d$ grows.

For $F_C(z)=e^{-a_C|z|^2}$ with $0<a_C\le1$, one may use
$2e^{-|\xi|}$ as a normalized Fourier envelope. For
$F_C(z)=(1+a|z|^2)^{-\nu}$, $\nu>1$, Gaussian integration gives
\[
 \frac{\widehat F_C(\xi)}{\int F_C}
 =\frac1{\Gamma(\nu-1)}\int_0^\infty
 t^{\nu-2}e^{-t-\pi^2|\xi|^2/(at)}\,dt
 \le2^{\nu-1}e^{-\sqrt2\pi|\xi|/\sqrt a}.
\]
The inequality follows from $t+L/t\ge t/2+\sqrt{2L}$.

\subsection{Compact and Complex-Pair Profiles}

Fix $0<\delta<1$, $p=2/(1+\delta)$, a positive radial $f_C$ on $\mathbb C$,
and a positive separately radial $g_D$ on $\mathbb C^2$. Here $p$ denotes
the Lebesgue exponent. Define
\[
 A_C=\int f_C^p,\quad I_C=\int f_C(z)f_C(z+1)\,dz,
 \qquad A_D=\int g_D^p,
\]
\[
 I_D=\int_{\mathbb R}\int_{\mathbb C^2}
 g_D(x,y)g_D(x+e^u,y+e^{-u})\,dx\,dy\,du,
\]
\[
 J_C=\log I_C-(1+\delta)\log A_C,
 \quad J_D=\log I_D-(1+\delta)\log A_D.
\]
Assume these integrals are finite and positive, the two endpoint log-energy
variables have finite second moments, the total energy is bounded below
and coercive, and \eqref{geo:pointwise-poisson} holds for the tensor
weight $F=\prod f_C^p\prod g_D^p$ uniformly in reciprocal deformations.
Section~\ref{pf:section} verifies these hypotheses for our profiles. They
give the following geometric criterion.

\begin{proposition}\label{geo:transfer}
Suppose the preceding arithmetic data exist with $d\to\infty$, and
$\log L_F(1)\le dC$. With $H,J,\ell$ fixed up to bounded total rounding
errors, set
\begin{equation}\label{geo:general-margin}
 \mathcal M(\theta)=J-\delta H-(1/2-\delta)\ell-C
 +(1-\delta)\log2+(1-\theta)\log\pi
 +(1-2\theta)J_C+\theta J_D.
\end{equation}
If this margin has a fixed positive lower bound over all occurring
signatures, there are arbitrarily large planar sets whose unordered
unit-distance count divided by $|U|^{1+\delta}$ tends to infinity.
\end{proposition}

\begin{proof}
Write $V_C=-\log f_C$, $V_D=-\log g_D$. The compact endpoint law has
density $f_C(z)f_C(z+1)/I_C$ with respect to $dz$. The pair endpoint law
has density $g_D(x,y)g_D(x+e^u,y+e^{-u})/I_D$ with respect to $dx\,dy\,du$.
These are probability laws. Reflection about the midpoint exchanges the
two endpoints without changing either law. Let the first-endpoint energy
means be $v_C,v_D$ and set
$v=(b v_C+c v_D)/d$. Take the common sublevel set
\[
 \Omega=\left\{\sum V_C+\sum V_D\le d(v+\eta)\right\},\qquad\eta>0.
\]
It is bounded. Its volume is at most
$A_C^bA_D^c e^{pd(v+\eta)}$, and the pointwise weighted bound gives
\[
 N_{\max}\le2e^{dB_\eta},\quad
 B_\eta=H+\log2-\ell+(1-2\theta)\log A_C+\theta\log A_D+pv+p\eta.
\]
Independence across blocks, endpoint exchange symmetry, Chebyshev's inequality and a
union bound put both endpoint energies within $\eta d$ of their means
with probability at least $1/2$ for large $d$, uniformly in signature.
Undoing the overlap density gives
$\int\Phi\ge\tfrac12 I_C^b I_D^c e^{2d(v-\eta)}$.
Equations \eqref{geo:mass-density} and \eqref{geo:joint-average}
give ordered average degree at least $\tfrac14e^{dA_\eta}$, where
\[
 \begin{aligned}
 A_\eta={}&J+\log2+(1-\theta)\log\pi-\ell/2-C
 +(1-2\theta)\log I_C+\theta\log I_D\\
 &-(1-2\theta)\log A_C-\theta\log A_D
 +(2-p)v-(2+p)\eta.
 \end{aligned}
\]
Since $p(1+\delta)=2$, we have
$A_\eta-\delta B_\eta=\mathcal M(\theta)-4\eta$.
The ordered ratio is at least
$2^{-2-\delta}\exp(d(\mathcal M-4\eta))$, apart from bounded rounding
factors. Choose $4\eta<\inf\mathcal M$.
A compact coordinate comes from a real place of $F$. Projection to that
embedding is injective on differences of affine ideal points and sends
all distinguished norm-one displacements to unit vectors. Division by two
gives the unordered conclusion. Since $u(U)\le |U|^2/2$ and
$0<\delta<1$, the ratio $u(U)/|U|^{1+\delta}$ can tend to infinity only
if $|U|\to\infty$. This gives a sequence of cardinalities, with no claim
at every sufficiently large cardinality.
\end{proof}
